\documentclass{article}
\usepackage[left=2cm, right=2cm, top=2cm, bottom=2cm]{geometry}
\usepackage{graphicx}
\usepackage{amsmath}
\usepackage{amssymb}
\usepackage{amsthm}
\usepackage{fancyhdr}
\usepackage{verbatim}
\usepackage{listings}
\usepackage{xcolor}
\usepackage{pdfpages}
\usepackage{cancel}
\usepackage{physics}
\usepackage{esint}
\usepackage{braket}

\lstset{
    language=C++,
    basicstyle=\ttfamily\footnotesize,
    keywordstyle=\color{blue},
    commentstyle=\color{green},
    stringstyle=\color{red},
    numbers=left,
    numberstyle=\tiny\color{gray},
    frame=single,
    breaklines=true,
    captionpos=b,
}


\title{MTH 446: Midterm \# 1 Study}
\author{Cliff Sun}

\newtheorem{theorem}{Theorem}[section]
\newtheorem{lemma}[theorem]{Lemma}
\newtheorem{definition}[theorem]{Definition}
\newtheorem{conjecture}[theorem]{Conjecture}
\newtheorem{proposition}[theorem]{Proposition}
\newtheorem{corollary}[theorem]{Corollary}
\newtheorem{one minute paper}[theorem]{One Minute Paper}

\pagestyle{fancy}
\lhead{\textbf{\thepage}\ \ \nouppercase{\rightmark}}
\chead{MTH 446: Midterm \# 1 Study}
\rhead{Cliff Sun}

\begin{document}

\maketitle

\section*{Lecture Span}
\begin{enumerate}
    \item Study for midterm 1
\end{enumerate}

\section*{Q1: General information about complex numbers}

\begin{enumerate}
    \item $|z|^2 = z\cdot\overline{z}$
    \item Normalize $z$ to find argument
    \item $\cos\theta$ and $\sin\theta$ values at various $\theta$'s 
\end{enumerate}

\section*{Q2: Roots of complex numbers}

Remember to divide $2\pi k$ by $n$, where $n$ is the root of the problem. Then, iterate through $k$. Also, $k=0$ is the principal root. 

\section*{Q3: Limits}

\begin{enumerate}
    \item Algebraic properties of limits and cont. functions
    \item All polynomials are continuous
    \item $\epsilon$-$\delta$ definition
\end{enumerate}

\section*{Q4: Demoire's formula}

\begin{equation}
    (\cos\theta + i\sin\theta)^n = \cos n\theta + i \sin n\theta
\end{equation}

\section*{Q5: Sets}

\begin{enumerate}
    \item Boundaries (for all $\epsilon > 0$, $D_{\epsilon}(z_0)$ intersects with interior $A$ and exterior $\mathbb{C}/A$)
    \item Interior, $D_{\epsilon}(z_0) \subset A$ for all $\epsilon$. 
    \item Exterior, the same but $\subset \mathbb{C}/A$
    \item Open sets don't contain their boundary, closed is the contrary.
    \item Closure 
    \item Accumulation points, deleted $\epsilon$ neighborhood. 
\end{enumerate}

\end{document}